\section{An exact Erd\H{o}s--Straus/S6 terminal-matrix crosswalk}
\label{sec:es-s6-crosswalk}

This section records an arithmetic relation suggested by comparing the
candidate's terminal presentation with the divisor-shell coordinates used in
the maintained Erd\H{o}s--Straus corpus~\cite{ClankersES2026}.  All maps below
are proved over the integers.  The conclusion is an exact change of relators;
it is not a geometric map from the reciprocal cubic surface to the candidate
complex threefold.

\subsection{The first-half shell and its divisor coordinates}

\begin{proposition}[first-half denominator]
\label{prop:es-first-half-denominator}
\Verified\ Let \(p\equiv1\pmod4\) be an odd prime.  If
\[
 \frac4p=\frac1x+\frac1y+\frac1z,
 \qquad x,y,z\in\mathbb Z_{>0},\qquad x\le y\le z,
\]
then
\[
 \frac p4<x\le\frac p2.
\]
\end{proposition}

\begin{proof}
The lower bound follows from \(1/x<4/p\).  Suppose \(x>p/2\).  Clearing
denominators and reducing modulo \(p\) gives
\[
 4xyz=p(xy+xz+yz),\qquad p\mid xyz.
\]
If \(p\mid x\), then \(x,y,z\ge p\), so the sum of the three reciprocals is
at most \(3/p\).  If \(p\mid y\), then \(y,z\ge p\), and
\[
 \frac1x+\frac1y+\frac1z<\frac2p+\frac2p=\frac4p.
\]
Both alternatives are impossible.  Hence \(z=kp\) for an integer \(k\ge1\),
and substitution gives
\[
 (4k-1)xy=kp(x+y).
\]
Neither \(x\) nor \(y\) is divisible by \(p\), so \(p\mid4k-1\).  Write
\(4k-1=\ell p\).  Since \(p\equiv1\pmod4\), one has
\(\ell\equiv3\pmod4\), hence \(\ell\ge3\).  Cancelling \(p\) gives
\(\ell xy=k(x+y)\).  Since \(x\le y\),
\[
 x\le\frac{2k}{\ell}
   =\frac{\ell p+1}{2\ell}
   =\frac p2+\frac1{2\ell}<\frac{p+1}{2}.
\]
But the integer \(x>p/2\) is at least \((p+1)/2\), a contradiction.
\end{proof}

Fix henceforth a prime \(p\equiv1\pmod{12}\), and put
\[
 L=\frac{p-1}{4},\qquad
 R_j=4j+3,\qquad
 a_j=\frac{p+R_j}{4},\qquad 0\le j<L.
\]
The preceding proposition shows that every solution's smallest denominator
lies in the exact candidate set \(\{a_j:0\le j<L\}\).  Notice that
\[
 4a_j=p+R_j,\qquad
 \gcd(R_j,p)=\gcd(R_j,a_j)=\gcd(p,a_j)=1.
\]

\begin{proposition}[complete two-channel divisor selector]
\label{prop:es-two-channel-selector}
\Verified\ Define
\[
 \begin{aligned}
 E_j&=\{u\in\mathbb Z_{>0}:u\mid a_j^2,\ R_j\mid4u+1\},\\
 M_j&=\{u\in\mathbb Z_{>0}:u\mid a_j^2,\ R_j\mid4u+p\}.
 \end{aligned}
\]
Then the Erd\H{o}s--Straus equation for \(p\) has a positive integral
solution if and only if at least one \(E_j\) or \(M_j\) is nonempty.
For \(u\in E_j\) set \(\varepsilon=0\), and for \(u\in M_j\) set
\(\varepsilon=1\).  With
\[
 T=\frac{4u+p^\varepsilon}{R_j},\qquad
 K=\frac{p^{1-\varepsilon}T+1}{4},
\]
the corresponding solution is
\[
 (x,y,z)=
 \left(a_j,\frac{p^\varepsilon a_jK}{u},pK\right),
\]
up to interchange of \(y\) and \(z\).
\end{proposition}

\begin{proof}
Fix \(x=a_j\), and abbreviate \(R=R_j\), \(a=a_j\), and \(N=pa\).
The remaining equation is
\[
 Ryz=N(y+z).
\]
It is equivalent to
\[
 (Ry-N)(Rz-N)=N^2.
\]
Both factors are positive.  Since \(\gcd(R,N)=1\), positive solutions are in
bijection with positive divisors \(d\mid N^2\) satisfying
\[
 d\equiv-N\pmod R;
\]
the complementary divisor \(e=N^2/d\) then also satisfies
\(e\equiv-N\pmod R\), and
\[
 y=\frac{N+d}{R},\qquad z=\frac{N+e}{R}.
\]

Because \(\gcd(p,a)=1\), the \(p\)-adic exponent of \(d\) is \(0\), \(1\),
or \(2\).  Complementation \(d\leftrightarrow e\) exchanges exponents \(0\)
and \(2\), so choose the exponent-two representative in that pair.  If
\(d=p^2u\), then \(u\mid a^2\), and
\[
 R\mid d+N
 \Longleftrightarrow R\mid pu+a
 \Longleftrightarrow R\mid4u+1,
\]
where \(4a\equiv p\pmod R\) and \(p\) is invertible modulo \(R\).  This is
the \(E_j\) channel.  If \(d=pu\), then
\[
 R\mid d+N
 \Longleftrightarrow R\mid u+a
 \Longleftrightarrow R\mid4u+p,
\]
which is the \(M_j\) channel.  This proves completeness of the two selector
sets.

For either channel, \(R_jT=4u+p^\varepsilon\).  Modulo \(4\), this says
\(-T\equiv1\), so \(T\equiv3\pmod4\); the displayed \(K\) is therefore a
positive integer.  Direct calculation using \(4a=p+R_j\) gives
\[
 R_jK=a+p^{1-\varepsilon}u,
 \qquad Ta=u+p^\varepsilon K.
\]
Since \(u\mid a^2\), the second identity implies
\(u\mid p^\varepsilon aK\).  Finally,
\[
 \frac{u}{p^\varepsilon aK}+\frac1{pK}
 =\frac{p^{1-\varepsilon}u+a}{paK}
 =\frac{R_j}{pa}
 =\frac4p-\frac1a,
\]
which proves the asserted formula and its converse.
\end{proof}

\subsection{The unimodular relation-matrix bijection}

Put
\[
 \lambda(\ell_0,\ell_1,\ell_2)
 =12\ell_0-4\ell_1-3\ell_2.
\]
By \cref{prop:terminal-group}, the candidate's ordered terminal relation
matrix in the ordered generator basis \((x,c)\), for
\(\ell=(\ell_0,\ell_1,\ell_2)\), is
\[
 \mathsf R(\ell)=
 \begin{pmatrix}
 3&-\ell_1\\
 4&\ell_2-4\ell_0
 \end{pmatrix},
 \qquad \det\mathsf R(\ell)=-\lambda(\ell).
\]

For \(\varepsilon\in\{0,1\}\), define the marked selector space
\[
\mathcal W_\varepsilon(p)=
 \left\{(j,u)\in\mathbb Z\times\mathbb Z_{>0}:
 \begin{array}{l}
 0\le j<L,\ u\in\mathbb Z_{>0},\ u\mid a_j^2,\\
 R_j\mid4u+p^\varepsilon
 \end{array}
 \right\}
\]
and the normalized divisor-supported terminal slice
\[
\mathcal H_\varepsilon(p)=
 \left\{(j,\ell_0,\ell_1)\in\mathbb Z^3:
 \begin{array}{l}
 0\le j<L,\ \lambda(\ell_0,\ell_1,1)=p^\varepsilon,\\
 T:=4\ell_0-1>0,\quad
 u:=\ell_1+jT>0,\quad u\mid a_j^2
 \end{array}
 \right\}.
\]

\begin{theorem}[exact marked selector/terminal-slice equivalence]
\label{thm:es-s6-selector-terminal-bijection}
\Verified\ The map
\[
 \Theta_{\varepsilon,p}:\mathcal W_\varepsilon(p)
 \longrightarrow\mathcal H_\varepsilon(p),
 \qquad
 (j,u)\longmapsto
 \left(j,\frac{T+1}{4},u-jT\right),
 \quad T=\frac{4u+p^\varepsilon}{R_j},
\]
is a bijection.  Under it, the two relation matrices are related by the
explicit unimodular row operation
\[
 \begin{pmatrix}1&-j\\0&1\end{pmatrix}
 \begin{pmatrix}R_j&-u\\4&-T\end{pmatrix}
 =
 \begin{pmatrix}3&-\ell_1\\4&1-4\ell_0\end{pmatrix}
 =\mathsf R(\ell_0,\ell_1,1),
\]
and both determinants equal \(-p^\varepsilon\).
\end{theorem}

\begin{proof}
For \((j,u)\in\mathcal W_\varepsilon(p)\), \(T\) is a positive integer and
the congruence computation in \cref{prop:es-two-channel-selector} gives
\(T\equiv3\pmod4\), so \(\ell_0=(T+1)/4\) is integral.  Set
\(\ell_1=u-jT\).  Since \(R_j=4j+3\),
\[
 \lambda(\ell_0,\ell_1,1)
 =3(T+1)-4(u-jT)-3
 =(4j+3)T-4u=p^\varepsilon.
\]
The positivity and divisor conditions in the definition of
\(\mathcal H_\varepsilon(p)\) follow from \(u=\ell_1+jT\).  Matrix
multiplication gives the displayed identity, and
\[
 \det\begin{pmatrix}R_j&-u\\4&-T\end{pmatrix}
 =-R_jT+4u=-p^\varepsilon.
\]

Conversely, from \((j,\ell_0,\ell_1)\in\mathcal H_\varepsilon(p)\) recover
\[
 T=4\ell_0-1,
 \qquad u=\ell_1+jT.
\]
Expanding \(\lambda(\ell_0,\ell_1,1)=p^\varepsilon\) gives
\(R_jT=4u+p^\varepsilon\).  Thus \(R_j\mid4u+p^\varepsilon\), and the
remaining positivity and divisor hypotheses are already part of
\(\mathcal H_\varepsilon(p)\).  These constructions are literal inverses.
\end{proof}

The normalization \(\ell_2=1\) automatically satisfies the candidate's odd
\(\ell_2\) admissibility condition.  It also forces \(3\nmid\ell_1\): if
\(3\mid\ell_1\), then \(\lambda(\ell_0,\ell_1,1)\) is divisible by \(3\),
whereas \(p^\varepsilon\equiv1\pmod3\).  Write
\(\operatorname{ES}(p)\) for the assertion that \(4/p\) is a sum of three
positive integral unit fractions.  Consequently
\[
 \operatorname{ES}(p)
 \Longleftrightarrow
 \mathcal H_0(p)\sqcup\mathcal H_1(p)\ne\varnothing.
\]
The \(E\)-channel is the slice with terminal parameter
\(p_{\mathrm{S6}}=1\); the \(M\)-channel is the slice with
\(p_{\mathrm{S6}}=p\).

Here the channel label is part of the datum.  The raw selector subsets need
not be disjoint: for \(p=13\), \(j=0\), and \(u=2\), one has
\(u\in E_0\cap M_0\).  Accordingly the occurrence statement above means the
tagged disjoint union
\(\{0\}\times\mathcal H_0(p)\sqcup\{1\}\times\mathcal H_1(p)\).
Moreover the map from tagged selector states to \emph{unordered} unit-fraction
solutions is not injective in the \(M\)-channel: complementation
\(u\leftrightarrow a_j^2/u\) interchanges \(y\) and \(z\).  For example,
\(p=13\), \(j=0\), and \(u=2,8\) give the same unordered solution
\(\{4,26,52\}\).  The theorem is a bijection with the marked terminal slice,
not a bijection with unordered solutions.

The tags are essential.  For every eligible prime \(q\), define the
fixed-prime forgetful map
\[
 \mathfrak f_q:\mathcal H_0(q)\longrightarrow\mathbb Z^3,
 \qquad
 \mathfrak f_q(j,\ell_0,\ell_1)=(\ell_0,\ell_1,1).
\]
For \(p=13\), the two \(E\)-channel selector states satisfy
\[
 \Theta_{0,13}(0,2)=(0,1,2),
 \qquad
 \Theta_{0,13}(1,5)=(1,1,2),
\]
so their marked images are distinct, while both are sent by
\(\mathfrak f_{13}\) to \((1,2,1)\).  Thus forgetting \(j\) destroys
injectivity.  To compare different prime fibres, define
\[
 \mathfrak F:
 \bigsqcup_{\substack{q\ \mathrm{prime}\\q\equiv1\pmod {12}}}
 \bigl(\{q\}\times\mathcal H_0(q)\bigr)
 \longrightarrow\mathbb Z^3,
 \qquad
 \mathfrak F(q;j,\ell_0,\ell_1)=(\ell_0,\ell_1,1).
\]
At \(q=73\), one has \(\Theta_{0,73}(1,5)=(1,1,2)\), so this tagged
image is also sent by \(\mathfrak F\) to \((1,2,1)\).  Hence forgetting
the external prime loses still more information.

\subsection{What the crosswalk proves, and what it does not}

The theorem identifies the complete fixed-\(p\), normalized, positively
oriented, \(j\)-marked divisor selector with an arithmetic slice of the
candidate's terminal relation matrices.  It explains a genuine structural
continuity between the earlier reciprocal-equation work and the terminal
algebra used here: the common object is a \(2\times2\) integral presentation,
and the passage between the two displayed matrices is an element of
\(\operatorname{SL}_2(\mathbb Z)\).

The orientation qualification is substantive.  The geometrically validated
candidate specialization used elsewhere in this workbench has terminal
parameter \(p_{\mathrm{S6}}=-1\), whereas the \(E\)-channel here has
\(p_{\mathrm{S6}}=+1\).  No result in this section identifies those two
oriented terminal data.  Thus the exact arithmetic equivalence does not, by
itself, place an Erd\H{o}s--Straus selector state on the particular complex
threefold reconstructed in the candidate-specific \(p=-1\) branch.

It does not prove that the selector is nonempty for every prime, so it does
not prove the Erd\H{o}s--Straus conjecture.  It also does not construct a map
from the reciprocal cubic surface to \(X\), identify selector divisors with
cycles or attaching maps in the candidate geometry, or transfer the
candidate's complex structure to the arithmetic surface.  Those would be
additional geometric assertions, not consequences of unimodular equivalence
of relation matrices.
